% 本文件由 LaTeX 排版 Agent 根据有效 translation.json 自动生成。
% author=Miscellaneous; work_id=0855; title=西满刻法在罗马的教导

\chapter{\Person{西满刻法}在\Place{罗马}的教导}

% structure: p0001 source=h1 target=omitted reason=page-title-already-rendered-by-parent

\Emph{\Person{西满刻法}在\Place{罗马}城的教导。}\par

在\Person{克劳狄凯撒}三年，\Person{西满刻法}离开\Place{安提约基雅}前往\Place{罗马}。当他前行时，他在\Emph{各个}国家宣讲我们主的话。当他快要到达那里时，许多人已听见\Emph{这件事}，就出去迎接他，整个教会都满心欢喜地接待了他。城里的一些首领，佩戴着帝国冠带的人，也来到他面前，为要见他并听他的话。当全城的人都聚集在他周围时，他站起来向他们讲话，向他们展示他的教义宣讲是怎样的。于是他开始这样对他们说：——\par

众人，罗马的百姓，全意大利的圣徒，请听我对你们说的话。今天我宣讲并宣告耶稣，天主子，祂从天降下，成为人，在我们中间如同\Emph{我们中的一员}，行了奇妙的异能、神迹和奇事，在我们面前，也在巴勒斯坦地所有的犹太人面前。你们自己也曾听闻祂所行的事：因为别的国家的人也因祂医治的名声和祂所施奇妙救助的传闻而来到祂面前；凡亲近祂的人，都因祂的话而得痊愈。正因为祂是天主，在医治的同时祂也赦免罪过：因为祂公开的医治为祂隐藏的赦免作证，证明那是真实可靠。因为这位耶稣，先知们曾在其奥秘言词中预报过，他们期待着看见祂并听到祂的话：祂从永远和亘古就与圣父同在；天主，在高处隐藏，在低处显现；荣耀的圣子，出自圣父，并与圣父、祂的圣神及祂统治的可怕权能同受荣耀。祂自愿被罪人之手钉在十字架上，并被提升到圣父那里，正如我和我的同伴所看见的。祂还要在自己的荣耀和祂圣天使的荣耀中再次来临，正如我们听见祂对我们所说的。因为我们不能说任何不是从祂那里听到的话，我们也不在祂的福音书上写任何祂自己没有对我们说的话：因为这言语的发出，是为了在众人要为闲话在审判之地交账的日子，堵住说谎者的口。\par

此外，因为我们原是渔夫，不精通典籍，因此祂也对我们说：\Quote{我要派遣圣神，就是\OriginalTerm{护慰者}{Paraclete}到你们那里，使祂教导你们所不知道的事；}因为我们所说的你们所听见的，乃是\Emph{祂}的恩赐。而且，藉着祂，我们也给病人带来帮助，给患病者带来痊愈：好使你们因听到祂的话并藉祂权能的帮助，而相信基督就是天主，天主子；并且脱离奴役的束缚，崇拜祂和圣父，并荣耀祂的圣神。因为当我们荣耀圣父时，我们也同时荣耀圣子；当我们崇拜圣子时，我们也同时崇拜圣父；当我们宣认圣神时，我们也宣认圣父和圣子：因为我们受命因父、及子、及圣神之名给相信的人授洗，使他们获得永生。\par

因此，要逃避这世界智慧的言语，其中毫无益处，而要亲近那些真实、可信、且在天主面前蒙悦纳的言语；它们的赏报已存留，它们的酬报\Emph{确凿}。如今，光明已升起在受造物之上，世界获得了心灵的眼睛，使每个人都能看见并明白：受造物不应代替造物主受崇拜，也不应与造物主一同受崇拜：因为凡是受造物，都是被造来崇拜其造物主的，不应像其造物主那样被崇拜。但这位来到我们中间的\Emph{那一位}，就是天主，天主子，按祂的本性，尽管祂将天主性与我们的人性结合，为的是藉祂的天主性更新我们的人性。因此，我们应当崇拜祂，因为祂应与圣父同受崇拜，而不应崇拜受造物，那些受造物是为崇拜造物主而被造的。因为祂本身就是真实与真理的天主；祂本身就是先于\Emph{一切}世界和受造物的那一位；祂本身就是真实的圣子，是来自崇高圣父的荣耀果实。\par

但是你们看到伴随并追随这些话的奇妙作为。无人会相信：看哪！时间短促自从祂升到圣父那里，看祂的福音如何飞越整个受造界——从而可以知晓并相信祂自己就是受造物的创造者，并且万物因祂的命令而存在。再者，你们看到太阳在祂死时变暗，你们自己也是见证。大地在祂被杀害时震动，幔子在祂死时裂开。关于这些事，总督\Person{比拉多}也是见证：因为他亲自派人将这些事奏报给\Person{凯撒}，这些事以及更多的事在他面前和你们城市的首领面前被宣读。为此，\Person{凯撒}对比拉多发怒，因为他曾不公正地听从了犹太人的劝说；因此他派人夺回了所授予他的权柄。这同一件事在罗马人的整个统治区域内公布并为人所知。因此，\Person{比拉多}所看见并奏报给\Person{凯撒}和你们尊贵的元老院的，正是我所宣讲和宣告的，我的同僚宗徒们也是如此。你们知道\Person{比拉多}不能向帝国政府报告未曾发生和他亲眼未见的事；但所发生并真实成就的——正是他所书写和报告的。此外，坟墓的看守也是那里所发生之事的见证：他们变得如同死人；当那些看守在\Person{比拉多}面前受审时，他们在他面前承认犹太人的大司祭给了他们多么大的贿赂，好叫他们可以说我们祂的门徒偷走了基督的尸体。看哪！你们已经听到了许多事；再者，如果你们不愿被所听见的事说服，至少被你们所看见的、因祂的名所行的奇事说服。\par

不要让术士\Person{西满}用非真实的表象欺骗你们，就是他向你们展示的，如同向没有悟性的人，他们不知道分辨所见所闻。所以，派人去把他带来，到全城聚集的地方，你们为我们选择一些标记，在我们面前行；你们看见谁行出同样的标记，你们就该相信谁。\par

他们立刻派人去把术士\Person{西满}带来；那些支持他意见的人对他说：\Quote{作为一个我们确信你里面有能力行任何事的人，请在我们众人面前行一个标记，让这位宣讲基督的加里肋亚人\Person{西满}也看见\Emph{它}。}当他们正这样对他说时，恰巧有一个死人经过，是他们中间一位首领和有名望之人的儿子。他们所有人聚集在一起，对他说：\Quote{你们中谁使这死人复活，谁就是真的，应当被相信和接纳，我们都要跟随他，听从他的一切话。}他们对术士\Person{西满}说：\Quote{因为你比加里肋亚人\Person{西满}先在这里，我们在他之前就认识你，请先展示你所拥有的能力。}\par

然后\Person{西满}不情愿地走近死人；他们把担架放在他面前；他向左右观看，又抬眼望天，说了许多话：有些大声说出，有些秘密地不出声。他拖延了很长时间，什么事也没有发生，什么也没有做成，死人依旧躺在担架上。\par

\Person{西满}\OriginalTerm{刻法}{Cephas}立刻勇敢地走向死人，在所有聚集的会众面前大声呼喊：\Quote{因耶稣基督之名，就是犹太人在\Place{耶路撒冷}钉死，我们所宣讲的那位，从那里起来。}\Person{西满}的话一说完，死人便活过来，从担架上站了起来。\par

全体人民看见并惊奇；他们对\Person{西满}说：\Quote{你所宣讲的基督是真实的。}许多人喊叫说：\Quote{让术士\Person{西满}，欺骗我们所有人的，被石头砸死。}但\Person{西满}，因为每个人都跑去看复活的死人，他从一条街到另一条街，从一家到另一家逃脱了，那天没有落入他们手中。\par

但全城的人拉住\Person{西满}\OriginalTerm{刻法}{Cephas}，欢喜并亲切地接待他；他不停息地以基督的名行奇迹和异事；许多人都信了他。此外，\Person{库普里努斯}，即那复生者的父亲，带\Person{西满}到他家里，合宜地款待他，他和他全家都相信基督是永生天主之子。许多犹太人和外邦人在那里成为门徒。当他的教导带来极大喜乐时，他在\Place{罗马}及周围的城市，以及\Place{意大利}人民的所有村庄建造了教会；他在那里以\Emph{\OriginalTerm{统治者的管理}{Superintendence of Rulers}之职}服事了二十五年。\par

过了这些年，\Person{尼禄凯撒}抓住他，把他关进监狱。他知道要钉死他；于是叫了执事\Person{安苏斯}，立他接替自己在\Place{罗马}作主教。这些事是\Person{西满}自己说的；此外，他也吩咐\Person{安苏斯}将他所\Emph{负责}的其他事教导众人，对他说：\Quote{除了\WorkTitle{新约}和\WorkTitle{旧约}，不可在众人面前读任何别的东西：这是不合规矩的。}\par

当\Person{凯撒}下令将\Person{西满}倒钉十字架，正如他自己向\Person{凯撒}所请求的，并砍下\Person{保禄}的头时，百姓中起了大骚动，整个教会中充满痛苦，因为他们失去了亲眼见到宗徒们的机会。向导\Person{伊苏斯}起来，夜间取走他们的尸体，以极大的荣誉埋葬了他们，那里成为许多人的聚集地。\par

就在那时，仿佛出于公义的审判，\Person{尼禄凯撒}放弃了他的帝国并逃跑，\Person{尼禄凯撒}对他们发动的迫害暂时停息。多年以后，宗徒们离开世界，得享大荣冠，当司铎的祝圣在\Place{罗马}全城和全\Place{意大利}进行时，\Place{罗马}城发生了大饥荒。\par

\Person{西满·刻法}的教导至此结束。\par

\SourceRecord{Translated by B.P. Pratten.；Ante-Nicene Fathers；Vol. 8.；Edited by Alexander Roberts, James Donaldson, and A. Cleveland Coxe.；Buffalo, NY: Christian Literature Publishing Co.,；1886.；<http://www.newadvent.org/fathers/0855.htm>.}
